\chapter{\textbf{Act II:} The Spectrum Decomposes The Activity Manifold}
\label{chapter:act-II}

\section{A critical scale, a mode, and a state matrix}
\label{sec:act2-handover}
\bigskip
% Outline item 1 (act2_manifold.md section 4). Carries no claim and points at no
% figure. One paragraph, no results: the section restates what arrives and poses
% the question. Sources: fact sheet 06. Three numerals only.
% FORBIDDEN here (sheet 06): any result; "the two axes" for sigma against f;
% "the biological cut" for f = 0; introducing f beyond the one clause below;
% "the gap keeps the Gram spectrum clear of the floor" as something Act I showed;
% "compact bulk".

% P1: the two objects received, the scope of the chapter, and the question.
Two objects arrive from Chapter~\ref{chapter:act-I}, and neither of them is a
statement about activity. The first is $\mathrm{sr\_crit}$, the nominal spectral
radius at which a substrate's own bulk edge reaches one. Because the gap ratio is
also the inverse bulk, every substrate carries its own, and the four the figures
draw do not carry the same one (Table~\ref{tab:act1-ladder}). That is what makes
the second matching axis available at all, and it is also why that axis is not
neutral: the substrates differ in one quantity only, so each axis holds fixed
exactly what the other lets vary. The second object is the Perron mode. The
previous chapter established that it stands clear of the bulk and that weight
placement is what puts it there, and it established nothing else about it, offering
no mechanism for why placement produces the gap and making no claim about capacity
or about any task. The mode arrives here as an object to be decomposed rather than
as an explanation, and that limit travels with it. Everything in this chapter sits
at the $N = 448$ % TIER0 3.12
parcellation and on non-negative weights, at $f = 0$; % SPINE, primary variable; A2 4.3.4
the sign fraction $f$ is an intervention rather than a second variable, and it
enters once, in Chapter~\ref{chapter:act-III}, where it is used to remove the
proposed cause.% TODO: repoint to the section label for 6.3 once chapter 6 exists.
The variable along which every claim below is made is the nominal spectral radius
$\sigma$, the operating point. What the reservoir's readout sees, however, is neither the
spectrum nor the operator that generates it. It is a state matrix: on the
memory-capacity task, $2500 \times 448$ % A2 2.5
real numbers, the states the reservoir passes through after its warmup is
discarded, which a ridge solver then inverts
(Chapter~\ref{chapter:methods}). So the question this chapter asks is which part
of the spectrum ends up where in the reservoir's state matrix.


\section{Three captures, and what each can be asked}
\label{sec:act2-captures}
\bigskip
% Outline item 2 (act2_manifold.md section 4). Methods. Carries no claim and
% points at no figure; it fixes definitions, operating points and scope.
% Sources: fact sheet 07, plus ACT2_CAPTURE_COVERAGE.md (tag % COV) and
% CONVENTIONS.md working rule 5 (tag % CONV).
% FORBIDDEN here (sheet 07): any ladder reading of the basis-alignment capture;
% "the manifold lives in low-frequency graph harmonics"; JUSTIFYING d_eff, which
% is section 5's; median-then-absolute-value on mean_state; "alpha was chosen at
% 1e-6 for d_eff"; "the Gram spectrum is a property of W"; "compact bulk".

% P1: the object, and the identity that poses the chapter's questions.
Everything in this chapter is measured on one object. Write the post-warmup states
as a matrix $A \in \mathbb{R}^{T_{\mathrm{eff}} \times N}$, one row per retained
timestep and one column per unit, assembled as Chapter~\ref{chapter:methods}
describes. What a linear readout can use is not $A$ itself but its second moment,
and there are two of those. The ridge solver inverts the un-centred Gram matrix
$G = A^{\top}A$; the covariance of the states about their own time-average is
formed instead from $\tilde{A} = A - \mathbf{1}m^{\top}$, where
$m = T_{\mathrm{eff}}^{-1}A^{\top}\mathbf{1}$ collects each unit's mean over time.
These are not the same object, and the difference between them is exact:
\begin{equation}
    A^{\top}A \;=\; \tilde{A}^{\top}\tilde{A} \;+\; T_{\mathrm{eff}}\,m\,m^{\top}.
    \label{eq:gram-split}
\end{equation}
Read the two terms on the right physically. The second is built from one number
per unit, the level that unit sits at on average, so it has \emph{rank one}: it
describes a single fixed pattern of activity that the network holds throughout the
run, and it says nothing about how the activity moves. The first is everything
that is left once that fixed pattern is subtracted, which is the part that varies
from one timestep to the next and therefore the only part that can carry
information about a changing input. Time-centring is the one subtraction that
separates them. The identity settles nothing on its own: it says the time-average
lives in a rank-one term, not that the Perron mode is that term, which is a
measurement and is made in Section~\ref{sec:act2-perron}.

% P2: the three questions the identity poses, and one capture per question.
Equation~\ref{eq:gram-split} poses three questions, and this chapter answers them
with three frozen captures, one apiece. \textbf{How large is the rank-one term?}
That is a question about a single scalar per run, and the
\emph{saturation-diagnostics} capture records it, along with the related quantities
describing where the units are sitting: how far the network has been driven into
the flat regions of its nonlinearity, and how strongly it has synchronised.
\textbf{What is the remaining, fluctuating term organised by?} A cloud of activity
in $N$ dimensions can be described in any orthonormal coordinate system, and asking
which one lines up with the fluctuations is asking which structural feature of the
network they follow. The \emph{basis-alignment} capture answers it by projecting
the time-centred states onto the leading $k$ vectors of three candidate bases and
recording the fraction of variance each holds: the eigenvectors of $W$, so the
dynamical modes themselves; the low-frequency harmonics of the graph Laplacian,
which are the smoothest patterns the wiring supports and are the network's
analogue of low-frequency Fourier modes \cite{atasoy2016human}; and a random orthonormal basis, which
supplies the chance level, since $k$ arbitrary directions out of $N$ capture about
$k/N$ of the variance by construction and a basis that fails to beat that has told
us nothing. \textbf{And how many directions of the fluctuating term does the
readout actually have available?} That is set by where the eigenvalues of $G$ sit
relative to the regularisation, so the \emph{Gram-spectra} capture stores the whole
spectrum, all $N$ eigenvalues per run, rather than any summary of it. Storing the
spectrum rather than the states is deliberate and is what makes the later sections
possible: a state matrix runs to tens of megabytes per cell where a spectrum is
$N$ floating-point numbers, % CONV working rule 5
and every dimensionality measure this chapter uses is a function of the spectrum
alone, so it can be recomputed afterwards at any regularisation strength without
re-running anything.

% P3: the scope limit on the narrowest capture, and the supercritical filter.
The basis-alignment capture is much the narrowest of the three, and its limits are
stated here rather than in a limitations section because they bind every sentence
drawn from it. It holds two substrates, the connectome and the degree-preserving
rewire, % COV 3.1
and each weight condition was captured at two spectral radii only: one low point
shared across conditions and one higher point of its own, which for the
all-positive weights used throughout this chapter is
3.0526. % COV 3.4
So nothing drawn from it may be read as a result about the null ladder, since it
carries neither an Erdős–Rényi rung nor the weight-permuted control, and the panels
of Figure~\ref{fig:F5} each sit at their own condition's operating point rather
than at one shared radius. Where anything below is called \emph{supercritical} it
means a nominal spectral radius of at least
3.05, % TIER0 3.2
the connectome's own critical scale of
3.078 % TIER0 2.1
rounded down to the nearest point of the capture grid; the threshold is structural
rather than tuned, and Section~\ref{sec:act2-counting} gives the evidence that it
is conservative.

% P4: the common-mode proxy and its sign convention. Stated once, travels everywhere.
The scalar that measures the rank-one term is the grand mean of the state matrix
over both time and units, $\bar{x} = N^{-1}\mathbf{1}^{\top}m$, which is simply the
mean of the same per-unit averages Equation~\ref{eq:gram-split} isolates; it is
large when every unit sits far from zero together, which is what a network
dominated by one common mode does. Its sign, however, is arbitrary, set by the
particular input realisation a seed draws, and measurably so:
48.8\% % COV 3.2
of cells in the saturation capture carry a negative value. The absolute value is
therefore taken per seed, \emph{before} any aggregation across seeds, and the order
is not a matter of taste. Since
$\bigl|\operatorname{med}_{s}(\bar{x}_{s})\bigr| \le
\operatorname{med}_{s}\bigl|\bar{x}_{s}\bigr|$ for any sample, taking the median
first can only understate the amplitude, and it understates it by an amount that
depends on how many of a substrate's seeds fall either side of zero, which differs
between substrates. The effect is a reordering rather than a common shift: in the
wrong order the connectome's value at $\sigma = 6$ falls from
0.759 % TIER0 3.7
to 0.638 % TIER0 3.12
while the weight-permuted null falls to
0.575, % A2 5.1
below it, reversing the comparison Section~\ref{sec:act2-perron} makes.

% P5: d_eff defined here, justified in section 5. Lag caveat in the footnote.
How many directions the readout has available needs a counting rule, and the
obvious ones do not survive contact with the data. The rank of $A$ is the textbook
answer and is useless in practice: the Gram spectra here decay smoothly over ten
orders of magnitude without ever reaching zero, so the numerical rank is $N$ for
every substrate and the question is not answered but hidden. Counting eigenvalues
above some threshold is no better, since the threshold is arbitrary and a
direction lying just either side of it is counted as one or as nothing. Both
attempts fail for the same reason: they treat availability as a property of the
state matrix alone, when what makes a direction usable is whether the
\emph{readout} can put weight on it. So rather than choose a rule, this thesis
reads one off the solver. A ridge fit already shrinks each direction by a factor
determined by how far its Gram eigenvalue stands above the regularisation, and
summing those factors gives the \emph{ridge effective rank}. With
$g_1 \ge \dots \ge g_N$ the eigenvalues of $G$ and $\alpha$ the ridge
regularisation strength,

\begin{equation}
    d_{\mathrm{eff}}(\alpha)
    \;=\; \sum_{i=1}^{N} \frac{g_i}{g_i + \alpha}
    \;=\; \operatorname{tr}\!\bigl[G(G + \alpha I)^{-1}\bigr].
    \label{eq:deff}
\end{equation}
Each direction contributes a factor between zero and one: near one when its Gram
eigenvalue stands clear of the regularisation floor, near zero when it lies well
below, and exactly one half at the crossover $g_i = \alpha$. The sum is therefore a
count of how many directions of the state matrix survive the regularisation,
softened at the boundary, and the regularisation acts as a floor because that is
what ridge regression does to a direction carrying little
variance.\footnote{The memory-capacity design is nested in the lag, so the Gram at
lag $k$ is the lag-zero Gram minus a positive semi-definite term of rank at most
$k$, while the capture stores the spectrum at lag zero. Weyl's inequality bounds
the difference by $0 \le d_{\mathrm{eff}}(0) - d_{\mathrm{eff}}(k) \le k$, since at
most $k$ directions can be displaced and each is worth at most one count. At the
deepest lag that permits a shift of up to 50 % A2 2.6
units; measured on real cells it is 0.31 to 0.82, % A2 2.6
a quarter of one per cent of the range reported across substrates in
Section~\ref{sec:act2-counting}, so the stored spectrum stands in for every lag and
the approximation cannot move an ordering.}
The trace form is why this is the natural count here rather than a descriptor
chosen to suit the argument: the ridge solution is
$\hat{w} = (G + \alpha I)^{-1}A^{\top}y$, so Equation~\ref{eq:deff} is the
effective degrees of freedom of the fitted readout in the standard
sense~\cite{hastie2009elements}. That has a consequence for $\alpha$ which is easy
to state and easy to get wrong. It is not chosen for the measure at all. It is the
readout's own, $10^{-6}$ % TIER0 3.3
on the memory-capacity task, taken from the same configuration the solver is handed
rather than from a second copy of the number that could drift. Why this count and
not a variance-weighted one is a result rather than a definition, and
Section~\ref{sec:act2-counting} makes that case; making it here would leave the
measure looking assumed.


\section{The Perron mode carries the mean}
\label{sec:act2-perron}
\bigskip
% Outline item 3 (act2_manifold.md section 4). Carries A2.1 and A2.2.
% Figure F4 (a, b). Chain step 2 in CROSS_ACT_SPINE.md: spectral gap to weak
% common-mode domination. Sources: fact sheet 08.
% FORBIDDEN here (sheet 08): "the connectome's Perron mode is small", or any
% reading in which weak domination means a weak leading eigenvalue; any ladder
% reading of panel (a); "the manifold lives in low-frequency graph harmonics";
% quoting the pooled 0.0001 as a memory-task number; "the leading mode carries
% none of the variance"; a signed median of the common-mode proxy; the
% common-mode amplitude quoted without its sigma; "hub inhibition collapses
% memory"; "compact bulk"; any graded curvature account.

% P1: what the identity leaves, and the two questions the two panels answer.
Equation~\ref{eq:gram-split} splits the readout's matrix into a rank-one term
carrying the time-average and a fluctuation term carrying everything that moves,
but it does not say which part of the spectrum ends up in which. That is the
question this section answers, and it has two halves. The first asks what the
fluctuations are organised by, once the time-average has been removed: if the
leading dynamical modes are still there, the decomposition is not clean and the
readout is working in their shadow. The second asks how large the rank-one term
gets, and how that depends on the substrate. Figure~\ref{fig:F4} answers them in
turn.

% P2: A2.1. Below chance, with the harmonics control and the aggregation.
Take the connectome at a supercritical operating point, remove each unit's mean
over time, and ask how much of what is left the leading eigenvectors of $W$
account for. The answer is that they account for less than nothing in particular.
Section~\ref{sec:act2-captures} fixed the scale this is read against: $k$ arbitrary
orthonormal directions out of $N$ hold about $k/N$ of the variance, so at $k = 1$
chance is
$1/448 = 0.0022$, % A2 2.1
and the measured random baseline of
0.0025 % A2 2.1
sits where it should. The top eigenmode of $W$ holds
0.0006, % A2 2.1
below its own chance level, and it is below it in
10 of 10 seeds at every $k \le 5$, % A2 F4
reaching chance only around
$k = 20$ % A2 F4
and tracking it thereafter. On its own that result is ambiguous, because it is also
what one would see if no structural basis whatever lined up with the fluctuations,
which would make the measurement uninformative rather than the modes unusual. The
graph harmonics settle it. On the same states, in the same panel, they sit
7.6 times % A2 F4
\emph{above} chance at $k = 1$, so the shortfall is specific to the modes the
Perron root belongs to and is not a property of structural bases in general. One
qualification travels with these numbers. The value of record for this comparison
is pooled across the capture's three tasks and reads
0.0001 % TIER0 3.12
against a baseline of
0.0023; % TIER0 3.12
the figures in this chapter are drawn on the memory-capacity task alone, where the
same quantity is several times larger, and the magnitude accordingly depends on
which is quoted. The \emph{direction} does not, which is what the claim rests on:
29 of the 30 % A2 2.1
task and seed combinations sit below their own random baseline.

% Figure F4.
\begin{figure}[ht]
    \centering
    \includegraphics[width=\textwidth]{F4.png}
    \caption[The Perron mode is a common mode]{%
    \textbf{The Perron mode is a common mode: it carries the time-average, and the
    fluctuations the readout uses are orthogonal to it.} Human $N = 448$ % TIER0 2.1
    self-built consensus, all-positive empirical weights, memory-capacity task,
    ten seeds, medians over seeds. \textbf{(a)} Fraction of the \emph{time-centred}
    state variance held by the leading $k$ vectors of three orthonormal bases, on
    the connectome alone at the supercritical operating point
    $\sigma = 3.05$; % COV 3.4
    both axes logarithmic, and not a statement about the null ladder. The dominant
    $W$ eigenmodes fall below the random baseline and reach it only around
    $k = 20$, % A2 F4
    while the graph harmonics on the same states sit well above it. The grey band
    is the across-seed range of the chance curve, which is the quantity the
    per-seed comparison is made against; a single random direction is far more
    variable. Pooled across the capture's three tasks the leading mode holds
    0.0001 % TIER0 3.12
    against a 0.0023 % TIER0 3.12
    baseline, the value of record. \textbf{(b)} Common-mode amplitude against
    nominal $\sigma$ for the four substrates, absolute-valued before medianing
    (Section~\ref{sec:act2-captures}). The connectome is the \textbf{least}
    dominated at every $\sigma$ above roughly 1.5,
    0.759 % TIER0 3.7
    at $\sigma = 6$ against 0.949 to 0.989, % TIER0 3.7
    despite carrying much the largest Perron root
    (Figure~\ref{fig:F2}). The nulls rise first because they become supercritical
    first; the dotted rule is the connectome's own critical scale and the shaded
    region lies above it.}
    \label{fig:F4}
\end{figure}

% P3: A2.2 and the inversion. Refers back to chapter 4 without re-quoting.
So the leading modes are absent from the fluctuations. Panel~(b) of the same
figure asks where they went, by measuring the common-mode amplitude across the
whole sweep in $\sigma$ for all four substrates. Every substrate's rises as it is
driven harder, which is what a network being pushed into the flat regions of its
nonlinearity does, and the nulls' rise first, since they reach their own critical
scales well before the connectome reaches its; those scales are
Chapter~\ref{chapter:act-I}'s and are not re-derived here. What matters is where
the four curves end up. At $\sigma = 6$ % TIER0 3.7
the connectome sits at 0.759 % TIER0 3.7
against the nulls' 0.949 to 0.989, % TIER0 3.7
and near criticality at $\sigma = 2$ % TIER0 3.7
the separation is wider still, 0.114 % TIER0 3.7
against 0.53 to 0.59. % TIER0 3.7
The connectome is the least common-mode dominated substrate on the ladder, and it
is the substrate carrying by far the largest Perron root. That inversion is the
finding of this section. It is not a paradox but a consequence of what
Chapter~\ref{chapter:act-I} measured: what separates the connectome is a gap
between its leading eigenvalue and its bulk, not a leading eigenvalue that is small
or a bulk that is unusual, and a gap of that kind lets the leading mode be driven
hard without the bulk following it. The nulls, whose bulk sits close behind their
Perron root, synchronise as a whole. Their units end up pinned near the extremes of
the nonlinearity together, which is what an amplitude approaching one means, and
almost nothing is left that varies.

% P4: what this establishes and what it does not; poses section 4.
Two things follow, and it is worth separating them from a third that does not. The
decomposition is clean: the leading mode carries the time-average and the bulk
carries what a readout could use, which is what licenses the rest of this chapter
to work in the fluctuation subspace and to treat the rank-one term as accounted
for. And the substrates differ in how much of their activity that rank-one term
consumes, in the direction opposite to the one the spectrum naively suggests. What
does \emph{not} follow is anything about capacity. Being least dominated by the
common mode is not the same as having the most directions a readout can act on, and
the step between them is the one this chapter has yet to measure; nothing above
touches how the fluctuation spectrum is arranged, only how much of the activity is
left in it. That arrangement is what the readout is finally sensitive to, through
the regularisation floor Equation~\ref{eq:deff} counts against, and it is where
Section~\ref{sec:act2-floor} goes next. Everything in this section, and everywhere
else in this chapter, is measured on non-negative weights; no sign manipulation
appears until Chapter~\ref{chapter:act-III}.





\section{The Gram spectrum against the ridge floor}
\label{sec:act2-floor}
\bigskip
% Outline item 4 (act2_manifold.md section 4). Carries A2.6 and A2.7.
% Figure F18 (a, b, c). Chain step 3 in CROSS_ACT_SPINE.md: the measured link
% from weak common-mode domination to usable dimensionality. Sources: sheet 09.
% FORBIDDEN here (sheet 09): "minimum" without "interior", in prose or caption;
% "1.8 to 3.9" as the span of the null minima; "raising alpha costs Erdos-Renyi
% little"; a low sensitivity read as good or bad on its own; "the gap puts the
% Gram spectrum clear of the floor" as though derived here; "the anisotropy
% hypothesis is confirmed", or the decay exponent used as support; "tracks",
% "predicts", or any regression language for the migration; "3.58 is 3.6";
% "four orders"; the published fractions quoted without the zero-strip; any
% sentence implying the four per-bin medians sum to 100; "compact bulk".

% P1: the question, the guard, and A2.6 stated as positions.
Section~\ref{sec:act2-perron} established that the connectome keeps more of its
activity in the fluctuation term than any of its nulls do. It said nothing about
how that activity is arranged, and arrangement is what the readout is sensitive
to: Equation~\ref{eq:deff} counts directions by how far their Gram eigenvalues
stand above the regularisation, so what matters is where each substrate's spectrum
sits relative to that floor. This is a measurement rather than a derivation.
Nothing here obtains the position of the Gram spectrum from the spectrum of $W$,
and the link back to the spectral gap runs through the common-mode account of the
previous section rather than around it. Measured supercritically, the positions
are strikingly different. The connectome holds
89.0\% % TIER0 3.6(e)
of its directions more than a decade clear of the floor, against Erdős–Rényi's
11.4\%; % TIER0 3.6(e)
read from the other end, 73.2\% % TIER0 3.6(c)
of Erdős–Rényi's spectrum lies at or more than a decade \emph{below} the floor
before the regularisation is raised at all, against the connectome's
5.1\%. % TIER0 3.6(c)
That is what $d_{\mathrm{eff}}$
412.9 % TIER0 3.6
against 74.8 % TIER0 3.6
counts, and the two statements are the same fact viewed from either side of the
floor.\footnote{All four statistics are computed on the design-Gram spectrum with
exact zeros removed first, as the implementation does it; the convention is
load-bearing rather than cosmetic, since on the full 448-direction denominator the
published fraction below the floor reproduces for no substrate (the connectome's
6.6\% % TIER0 3.6
reads 7.03\% % TIER0 3.6(a)
and Erdős–Rényi's 79.4\% % TIER0 3.6
reads 83.59\%), % TIER0 3.6(a)
and the error grows down the ladder because the count of numerically rank-deficient
directions does. The other three statistics are untouched: a zero contributes
exactly nothing to $d_{\mathrm{eff}}$ and to the sensitivity below, and is never
within a decade of the floor. Figure~\ref{fig:F18}(a) shows the zeros as a bin of
their own rather than folding them into the fraction below the floor, which makes
the exclusion visible rather than merely stated.}
These are per-bin medians over the fifty supercritical cells of each substrate, so
the bins partition each individual cell exactly while four medians taken separately
are not constrained to sum to one hundred, and the connectome's come to
102.1\%. % TIER0 3.6(e)
Panel~(a) of Figure~\ref{fig:F18} therefore groups its bars rather than stacking
them, since a stacked bar would assert a partition the medians do not form.

% Figure F18.
\begin{figure}[ht]
    \centering
    \includegraphics[width=\textwidth]{F18.png}
    \caption[Where the Gram spectrum sits relative to the readout floor]{%
    \textbf{Where the Gram spectrum sits relative to the readout floor is what
    sets usable dimensionality.} Memory-capacity task, all-positive weights,
    $N = 448$, % TIER0 2.1
    ten seeds. \textbf{(a)} Design-Gram eigenvalues by position relative to the
    floor, supercritically, as a percentage of all 448 directions; medians over
    the fifty cells per substrate, with $d_{\mathrm{eff}}$ carried in the legend
    because it is a property of the whole spectrum rather than of any one bin.
    Exact zeros form their own bin and are excluded from the fraction below the
    floor. Bars are grouped, not stacked: the bins partition each individual cell
    exactly, but four medians taken separately need not sum to one hundred.
    \textbf{(b)} Sensitivity of $d_{\mathrm{eff}}$ to the floor against spectral
    radius, whole curve drawn. It vanishes at both ends, so the minima that matter
    are the \emph{interior} dips, marked per substrate; the shaded low-radius limb
    is where the reservoir is too weakly driven to have much spectrum above the
    floor at all, and is excluded from the search. \textbf{(c)} The measured
    ridge-optimal radius against regularisation strength, which migrates toward
    each substrate's own interior dip. The three nulls coincide exactly at every
    value, so their markers nest rather than being offset. Panels~(a) and~(b) use
    the probe radius grid and~(c) the coarser sweep grid, so
    3.58 % TIER0 3.6(b)
    and 3.6 % TIER0 3.6
    are two grids' nearest points to one place rather than the same number. The
    correspondence in~(c) is consistent-with; nothing is fitted.}
    \label{fig:F18}
\end{figure}

% P2: the rate, with the both-ends warning taken off the algebra first.
How much a substrate stands to lose if the regularisation is raised is a different
question from where its spectrum sits, and the two are easy to conflate.
Differentiating Equation~\ref{eq:deff} gives the rate directly:
\begin{equation}
    -\frac{\mathrm{d}\,d_{\mathrm{eff}}}{\mathrm{d}\log\alpha}
    \;=\; \sum_{i=1}^{N} \frac{g_i\,\alpha}{(g_i + \alpha)^{2}},
    \label{eq:floor-sensitivity}
\end{equation}
in units of directions lost per e-fold of $\alpha$. Each term peaks, at one
quarter, when a direction sits exactly \emph{at} the floor, and falls away as it
moves in either direction. So the sum is small for a substrate whose spectrum
stands well clear of the floor and equally small for one whose spectrum has already
fallen far below it, and a low value never interprets itself: only the positions
above tell the two cases apart. Erdős–Rényi is the second case, and reading its
low absolute sensitivity as robustness would invert the truth. Per unit of
surviving dimensionality it is the \emph{most} floor-sensitive substrate on the
ladder, losing 29.1\% % TIER0 3.6(c)
of its effective rank per decade of $\alpha$ against the connectome's
5.6\%. % TIER0 3.6(c)
The absolute losses are nearly the same at the two ends of the ladder,
21.8 % TIER0 3.6(c)
directions against 23.3; % TIER0 3.6(c)
the connectome simply has some four hundred directions to lose them from and
Erdős–Rényi has around seventy-five.

% P3: A2.7 first half, the interior minima.
Neither the positions nor the rate is a fixed property of a substrate, because both
depend on where the substrate is being operated. Panel~(b) draws
Equation~\ref{eq:floor-sensitivity} across the whole sweep, and the curve is
two-humped: it falls to zero at zero radius, where a reservoir barely driven at all
has almost no spectrum above the floor to lose, rises as directions are pushed up
through the floor, and falls again once they are clear of it. The minima that
matter are therefore the \emph{interior} dips between the two humps, and they sit
at different radii for different substrates: the connectome's near
3.58, % TIER0 3.6(b)
every null's between 1.58 and 2.00. % TIER0 3.6(b)
The separation is not an artefact of taking medians, holding in
8 of 10 % TIER0 3.6(b)
seeds for the connectome and 10 of 10 % TIER0 3.6(b)
for each null. The qualifier is not droppable: the connectome's lowest sensitivity
anywhere over non-zero radius is not its interior dip at all but a value of
2.687 % TIER0 3.6(b)
out at radius 0.4211, % TIER0 3.6(b)
below the dip's 5.785, % TIER0 3.6(b)
and read as a global minimum this half of the result does not reproduce. That is
the same trap as the rate: a low sensitivity out on the left-hand limb means a
reservoir with nothing to lose.

% P4: A2.7 second half, the migration, with the retraction folded in.
This radius-dependence closes a loose end. Raising the regularisation penalises
operating points where many directions sit at the floor, so each substrate's best
operating point should move toward its own interior dip, and the measured optima do
exactly that. Across a regularisation grid spanning five orders of magnitude in
four steps, % TIER0 3.6(d)
the connectome's ridge-optimal radius moves from
2.4 to 3.6 % TIER0 3.6
while every null moves once, from 1.2 to 1.6, % TIER0 3.6
each toward its own dip. Two caveats travel with that and are not detachable from
it. The two halves sit on different radius grids, so
3.58 and 3.6 % TIER0 3.6
are two grids' nearest points to one place rather than one number recovered twice.
And the correspondence is consistent-with rather than fitted: what is shown is that
the direction and the endpoint agree for all four substrates, and nothing is
regressed on anything. It is worth recording what this replaces, since an earlier
account explained the same migration differently. The proposal was that the
connectome's optimum moves because its covariance is more anisotropic, so raising
the regularisation strips its directions faster. That is rejected on its own
measurement: of the three substrates tested the connectome has the
\emph{shallowest} covariance decay, so by that measure it is the least anisotropic
rather than the most, and the participation ratio is flat across them at
1.21 to 1.29. % TIER0 3.6
The rejection stands. What survives is the account above, refitted at the other end
of the spectrum, the decay exponent having been fitted over the top decile while
what governs a ridge readout is the bottom.

% P5: handover into section 5.
One consequence of these positions deserves stating before the next section uses
it. The floor sits some ten orders of magnitude % A2 4.5.3
below the leading variance, so a substrate can have hundreds of directions
standing clear of it that carry almost none of the variance between them. Whether
those directions are counted at all depends entirely on how a dimensionality
measure weights them, and the measure the field reaches for first weights by
variance.




\section{Which counting scheme sees it}
\label{sec:act2-counting}
\bigskip
% Outline item 5 (act2_manifold.md section 4). Carries A2.4 and A2.5, together
% contribution 6. Figure F6 (a, b, c). Chain step 4 in CROSS_ACT_SPINE.md.
% Sources: fact sheet 10.
% FORBIDDEN here (sheet 10): any sentence carrying one number from the seven-rung
% unit and one from the 350-cell unit; the rung-index correlation quoted as the
% result; "+0.54" for its PR value (sign corrected to negative); "PR is wrong";
% "d_eff is a better measure in general"; PR's failure as a separate finding; any
% mechanism for why the low-variance directions carry memory; 431 and 413 used
% interchangeably; "peak d_eff shows the connectome uses more directions";
% "the manifold lives in low-frequency graph harmonics"; "compact bulk".

% P1: the range first, because the range is the mechanism.
Put the same seven substrates through the same task and measure the dimensionality
of the states they produce, and the answer depends on which measure is used to an
extent that ought to be alarming. The ridge effective rank moves
5.5-fold % TIER0 3.12
across them, from 75 to 413 % TIER0 3.12
of the 448 directions available. The participation ratio, the variance-weighted
count the field reaches for first, moves
16\%, % TIER0 3.12
from 1.19 to 1.38. % TIER0 3.12
These are not different experiments. They are two functions of the same stored
spectra, computed on the same cells, and one of them says these substrates differ
enormously while the other says they are barely distinguishable. Measured memory
capacity moves roughly three-fold % A2 2.2
over the same substrates, so there is something there to be seen. The question is
what makes one measure blind to it.

% P2: both are weighted counts over one matrix; the mechanism in one cell.
The two measures are more alike than their disagreement suggests, which is what
makes the disagreement informative. Both are exactly sums of one weight per
direction of the same state matrix. For the ridge effective rank that is
Equation~\ref{eq:deff} as written, each direction weighted by
$g_i/(g_i + \alpha)$, how far its Gram eigenvalue clears the regularisation floor.
For the participation ratio it is true but less obvious. Writing
$p_i = \lambda_i / \sum_j \lambda_j$ for direction $i$'s share of the state
covariance,
\begin{equation}
    \mathrm{PR}
    \;=\; \frac{\bigl(\sum_i \lambda_i\bigr)^{2}}{\sum_i \lambda_i^{2}}
    \;=\; \sum_{i=1}^{N} \frac{p_i}{\sum_j p_j^{2}},
    \label{eq:pr}
\end{equation}
so it too is a sum of per-direction weights, each direction counted by its share of
the variance. Both are therefore counts, and each curve's area in
Figure~\ref{fig:F6}(a) is the number it is labelled with. What the panel shows on
one connectome cell is how differently the two weight the same directions: the
ridge weight is above one half for
438 of the 448 % A2 2.4
directions, while the variance weight has collapsed by the fifth, with
two directions carrying 95\% % A2 2.4
of a participation ratio of 1.28. % A2 2.4
The covariance and the Gram are formed on one and the same post-warmup state
matrix, so the difference between them is one of weighting and not of data.

% Figure F6.
\begin{figure}[ht]
    \centering
    \includegraphics[width=\textwidth]{F6.png}
    \caption[Memory lives in low-variance directions that a variance-weighted
    measure discounts]{%
    \textbf{Memory lives in hundreds of low-variance directions, which a
    variance-weighted measure discounts.} Memory-capacity task, all-positive
    weights, $N = 448$, % TIER0 2.1
    supercritical cells. \textbf{(a)} The per-direction weight each measure
    assigns, on \emph{one} connectome cell, with a zoom over the first twelve
    directions on the same linear scale. Both measures are exact sums of these
    weights, so each shaded area is the number beside it. This cell's effective
    rank is 431; % A2 2.4
    the connectome's median over all fifty supercritical cells, plotted in~(b), is
    413. % TIER0 3.6
    \textbf{(b, c)} The consequence across the seven substrates of the wider null
    ladder, medians over fifty cells each, against \textbf{measured} memory
    capacity. Both axes are zero-based, so the spread of the points is comparable
    between panels: over the same substrates the effective rank moves
    5.5-fold % TIER0 3.12
    while the participation ratio moves 16\%. % TIER0 3.12
    Panel~(b)'s upper limit is the hard ceiling at $N$. Both correlations reported
    in the text are against measured memory capacity, never against ladder
    position, which is a different quantity and is not the claim.}
    \label{fig:F6}
\end{figure}

% P3: PR's failure is a consequence of section 4.
That collapse is not a separate finding. It follows from where the spectrum sits,
which Section~\ref{sec:act2-floor} has just measured. The regularisation floor lies
some ten orders of magnitude below the leading variance, so a direction can stand
far clear of the floor, and be fully available to the readout, while holding a
share of the variance too small to register in a weighted count. Hundreds of the
connectome's directions are in exactly that position. A variance-weighted measure
does not fail here through any defect of its own; it fails because the quantity it
weights by and the quantity that determines availability have come apart, and they
come apart precisely where the previous section found the substrates to differ.

% P4: the test against measured MC, one aggregation unit per sentence.
Whether that matters is a question the data can answer, by asking which measure
tracks what the readout actually achieves. Memory capacity enters here as the
quantity a dimensionality measure is being judged against, not as a performance
claim about any substrate: what is being ordered is the counting schemes, not the
connectome and its nulls. Across the seven substrates, taking one median per
substrate, the ridge effective rank orders them against measured capacity at a rank
correlation of +1.000 % TIER0 3.12
and the participation ratio at +0.107. % TIER0 3.12
Pooled instead over all 350 % TIER0 3.12
individual cells within the supercritical regime, the pair reads
+0.998 % TIER0 3.12
against +0.308. % TIER0 3.12
Those are two separate tests on two aggregation units, and a number from one may
not be set against a number from the other, though the two agree on everything that
matters. Correlating either measure against position in the null ladder is a
different quantity again, weak and negative for both, and is not the claim being
made. One limit belongs with this, since it bounds what the ordering can mean:
every substrate's effective rank peaks within a few per cent of the hard ceiling at
$N$, % TIER0 3.2
so peak dimensionality is unresolvable at this parcellation and the ordering above
lives in the decay region beyond the peak rather than at it.

% P5: observer-relativity owned, then scope.
The natural way to write this result is that the participation ratio is wrong, and
that would be the wrong way. It measures what it claims to measure, which is how
concentrated the variance of the states is, and for a question about the shape of
the activity cloud it is the right instrument. The claim here is narrower and, I
think, more interesting: that quantity is not the one a ridge readout can act on,
so a measure of dimensionality is only meaningful relative to what is reading the
system out. Two measures of the same states can disagree by 5.5-fold against 16\%
without either being in error, and the choice between them is a question about the
observer rather than about the data. The demonstration is empirical and its scope
should be read narrowly: one task, one family of substrates, one readout. That
total capacity is bounded by the number of linearly independent state variables is
established, and this is a case of that bound rather than a rival to
it~\cite{dambre2012information}; and the effective rank counted here is a property
of the states a task drives the network through, not of the recurrent matrix, which
is a distinct quantity that shares the
name~\cite{clark2025connectivity}.




\section{Sign selects the basis}
\label{sec:act2-sign}
\bigskip
% Outline item 6 (act2_manifold.md section 4). Carries A2.3, the one claim in the
% act that is deliberately NOT a contribution. Figure F5 (a, b, c). Short, and
% framed as a statement about the decomposition rather than about sign.
% Sources: fact sheet 11.
% FORBIDDEN here (sheet 11): "the manifold lives in low-frequency graph
% harmonics"; any ladder reading of this capture (two substrates only);
% "sign-gating of the manifold transition is a new result"; "balancing the signs
% improves the representation"; quoting the capture values as evidence of how
% much structure the bases hold; treating this as a result about the sign
% fraction f, which is chapter 6's intervention; "compact bulk"; any graded
% curvature account.

% P1: the question the decomposition leaves open, and why it is settled here.
One question has been left open since Section~\ref{sec:act2-perron}, and it is
worth returning to now that the previous section has supplied what is needed to
answer it properly. With non-negative weights the leading eigenmodes of $W$ were
absorbed into the rank-one term of Equation~\ref{eq:gram-split}, which leaves the
fluctuations to be organised by something other than the dynamical modes
themselves. What organised them, in Figure~\ref{fig:F4}(a), was the low-frequency
harmonics of the graph Laplacian, the smoothest patterns the wiring supports, which
enter the comparison because connectome-defined harmonic modes have been shown to
structure measured cortical
activity~\cite{atasoy2016human}. That result is easy to read as a fact about the
connectome. It is not. It is a fact about non-negativity, and the way to show as
much is to remove the non-negativity and leave everything else where it is.
Balancing the signs of the weights destroys the Perron guarantee, so there is no
dominant common mode for the leading eigenmodes to be absorbed into, and the
prediction is simply that they should come back.

% P2: A2.3, stated as the swap, with the qualifier that stops the overstatement.
They do. Figure~\ref{fig:F5} repeats the basis comparison across three weight
conditions on one substrate and one task: the empirical weights used throughout
this chapter, the same magnitudes with their signs balanced, and gaussian weights.
At $k = 10$ the harmonics hold 0.168 % TIER0 3.12
of the time-centred variance against the $W$ eigenmodes'
0.004 % TIER0 3.12
with non-negative weights; balancing the signs reverses it to
0.019 against 0.081; % TIER0 3.12
and with gaussian weights the eigenmodes hold
0.522 % TIER0 3.12
against 0.012. % TIER0 3.12
The ordering swaps, it swaps in the same direction on the other two tasks the
capture holds, and the topology never changes: the same graph, the same edges, the
same weight magnitudes in the first two conditions, with only the signs
redistributed. What is claimed is the swap. It is worth being explicit that nothing
stronger is available, because the numbers invite it: on the all-positive substrate
neither basis holds much of anything, the whole range across bases and tasks
running 0.04 to 0.17, % TIER0 3.12
so a claim that the fluctuations live in graph harmonics would overstate what has
been measured. Which basis leads is the measurement; how much either one captures
is a separate question, and the answer to that one is "not much".

% Figure F5.
\begin{figure}[ht]
    \centering
    \includegraphics[width=\textwidth]{F5.png}
    \caption[Sign composition selects which structural basis the fluctuations
    occupy]{%
    \textbf{Which structural basis the fluctuations occupy is set by the weight
    signs, not by the topology.} Fraction of time-centred state variance held by
    the leading $k$ vectors of each basis, connectome substrate, Lorenz task,
    $N = 448$, % TIER0 2.1
    medians over ten seeds, each panel at its own condition's supercritical
    operating point. \textbf{(a)} With \textbf{non-negative} weights the graph
    harmonics lead throughout and the dominant $W$ eigenmodes run below chance.
    \textbf{(b)} With the \textbf{same weight magnitudes} and balanced signs the
    ordering has already reversed. \textbf{(c)} With \textbf{gaussian} weights the
    eigenmodes dominate outright. The rule marks $k = 10$, where the values quoted
    in the text are read; the swap is visible across the whole curve, and the
    topology is identical in all three panels. On the all-positive substrate
    neither basis holds much in absolute terms, which is why what is claimed is the
    ordering. Two substrates only: this capture cannot speak to the null ladder.}
    \label{fig:F5}
\end{figure}

% P3: the agreement with section 5, then priority and the forward clause.
The low capture in panel~(a) is not a disappointment but a confirmation, and it is
here that the two halves of this chapter meet without having been made to.
Section~\ref{sec:act2-counting} found that on this substrate, supercritically, the
readout has around 413 % TIER0 3.6
of 448 directions available to it. Fluctuations spread over hundreds of directions
are fluctuations that no ten-vector basis should be able to hold, so the small
captures in panel~(a) are what the other measurement predicts. The relationship
holds at the other end too: where a condition's effective rank is low the capture
is high, reaching 0.886 % TIER0 3.12
on the gaussian weights for one of the tasks this capture covers. Two probes
measuring different things on different bases agree about which substrates and
conditions spread their activity widely, and neither was fitted to the other. Two
final points, both of which place this section rather than extending it. The first
is priority: that weight sign composition gates this transition is largely
pre-empted by Krauss et al.~\cite{krauss2019weight}, and it is reported here as
confirmatory. It earns its place in this chapter because the counting argument of
the previous section needs the basis established first, not because it is new. The
second is scope. The signed and gaussian conditions above are weight families
compared side by side, nothing has been done to the connectome to produce them, and
the question of what happens when negative weights are introduced as an
intervention on the real substrate belongs to Chapter~\ref{chapter:act-III}, where
it is used to remove the common mode and check whether its consequences go with it.








